% !PRINTED 452
\weberpage{452}
%<<CONT
one obtains the field of complex numbers $x+yi$, with rational $x, y$, which we shall call $\mathfrak{J}$. By the adjunction of a variable $u$ to the field $\mathfrak{R}$ one obtains the field which consists of all integral and fractional rational functions of $u$ with rational coefficients; by the adjunction of $u$ and $i$ to $\mathfrak{R}$, the field of the rational functions of $u$ whose coefficients are numbers in $\mathfrak{J}$, and so on.

Let it finally be remarked that a \gesperrt{field} of quantities is often also called, following \gesperrt{Kronecker}'s example, a \gesperrt{domain of rationality}. It is sometimes useful to employ the expression domain of rationality for a field when it is to be expressed thereby that, in a particular problem, the quantities of this field are to be regarded as known or rational.

\begin{center}
\S.\,141.

\gesperrt{Functions in a Field.}
\end{center}

We often have occasion to consider integral rational functions of a variable $x$
\weq{1}{f(x) = a_0 x^m + a_1 x^{m-1} + \cdots + a_{m-1} x + a_m}
whose coefficients belong to a particular field $\Omega$, and which we shall call \gesperrt{functions contained in} $\Omega$, or briefly \gesperrt{functions in} $\Omega$. They are not to be confused with functions which, say, \gesperrt{belong} to the field $\Omega$, when $\Omega$ is a function field.

In the third section we saw that, when the coefficients $a_0, a_1 \ldots a_m$ have definite numerical values, there exist $m$ roots $\alpha_1, \alpha_2 \ldots \alpha_m$ of $f(x)$, which likewise have definite numerical values, among which, however, several equal to one another occur when the discriminant $\Delta(a)$ of the function (1) vanishes. If the coefficients $a$ of $f(x)$ are independent variables, then for every definite system of values of the $a$ there is a definite system of values of the $\alpha$, and the $\alpha$ vary continuously with the $a$ (\S.\,40). So long, therefore, as we restrict the variation of the $a$ to a sufficiently narrow region, in which no such values occur for which the discriminant $\Delta(a)$ vanishes, the $m$ quantities $\alpha$, as functions of the $a$, are completely determined and distinguished from one another as soon as the value of each of them is given for some particular system of values of the $a$.
%CONT>>
